\documentclass[11pt,a4paper]{article}
\usepackage[utf8]{inputenc}
\usepackage[T1]{fontenc}
\usepackage[catalan]{babel}
\usepackage{lmodern}
\usepackage[sfdefault,lining]{FiraSans}
\usepackage[a4paper,top=2.3cm,bottom=2.2cm,left=1.8cm,right=1.8cm,headheight=1.3cm,headsep=0.4cm,footskip=1cm]{geometry}
\usepackage{amsmath,amssymb}
\usepackage{icomma}
\usepackage{graphicx}
\usepackage[table]{xcolor}
\usepackage{tikz}
\usepackage{fontawesome5}
\usepackage{ccicons}
\usepackage{fancyhdr}
\usepackage{lastpage}
\usepackage{tasks}
\usepackage{needspace}
\usepackage[most]{tcolorbox}
\usepackage{array}
\usepackage[hidelinks,pdftitle={Solucions de la fitxa de la unitat 8: Límits i continuïtat. 1r de Batxillerat MI},pdfsubject={Matemàtiques I},pdfkeywords={límits, límits a l'infinit, comparació d'infinits, indeterminacions, límits laterals, continuïtat, tipus de discontinuïtat, funcions a trossos, paràmetres, asímptotes, branques parabòliques}]{hyperref}

\definecolor{navy}{HTML}{1B2748}
\definecolor{accent}{HTML}{E8603A}
\definecolor{teal}{HTML}{1F9E90}
\definecolor{soft}{HTML}{EEF1F8}
\definecolor{muted}{HTML}{5D6685}

\setlength{\parindent}{0pt}
\setlength{\parskip}{3pt}
\settasks{label=\alph*), label-format=\bfseries\color{navy}, label-width=1.3em, item-indent=1.8em, column-sep=0.6em, after-item-skip=5pt, before-skip=3pt, after-skip=4pt}

% ---------- capçalera i peu ----------
\newcommand{\llicencia}{\href{https://creativecommons.org/licenses/by-nc-sa/4.0/deed.ca}{\ccbyncsa}\ \ Llicència Creative Commons Reconeixement-NoComercial-CompartirIgual 4.0 Internacional (CC BY-NC-SA 4.0)}
\pagestyle{fancy}
\fancyhf{}
\fancyhead[L]{\raisebox{-0.35\height}{\includegraphics[height=0.95cm]{logo-gva.png}}}
\fancyhead[C]{\small\color{muted}Solucions de la fitxa · Unitat 8: Límits i continuïtat · 1r Batxillerat MI}
\fancyhead[R]{\raisebox{-0.35\height}{\includegraphics[height=0.95cm]{logo-vermellar.png}}}
\fancyfoot[L]{\scriptsize\color{muted}\llicencia}
\fancyfoot[R]{\small\color{muted}\thepage}
\renewcommand{\headrulewidth}{0.4pt}
\renewcommand{\headrule}{\color{muted!40}\hrule height \headrulewidth}
\fancypagestyle{primera}{\fancyhf{}\renewcommand{\headrulewidth}{0pt}
  \fancyfoot[L]{\scriptsize\color{muted}\llicencia}
  \fancyfoot[R]{\small\color{muted}\thepage}}

% ---------- nivells i exercicis ----------
\newcommand{\estrelles}[1]{{\color{accent}\ifcase#1\or\faStar\faStar[regular]\faStar[regular]\or\faStar\faStar\faStar[regular]\or\faStar\faStar\faStar\fi}}
\newcounter{exer}
\newcommand{\ex}[1]{\par\needspace{4\baselineskip}\medskip\refstepcounter{exer}\noindent{\color{navy}\bfseries\theexer.}\ \ }

\newcommand{\seccio}[2]{\par\needspace{16\baselineskip}\bigskip
  \begin{tcolorbox}[colback=navy,colframe=navy,coltext=white,arc=2mm,boxrule=0pt,left=3mm,right=3mm,top=1.2mm,bottom=1.2mm]
  {\large\bfseries #1\quad #2}
  \end{tcolorbox}\vspace{-2pt}}

\newcommand{\exf}[1]{\par\medskip\refstepcounter{exer}\noindent{\color{navy}\bfseries\theexer.}\ \ }
\providecommand{\gt}{>}\providecommand{\lt}{<}
% quadre de recordatori (només a la fitxa d'exercicis)
\newcommand{\recorda}[1]{\par\smallskip\begin{tcolorbox}[colback=accent!7,colframe=accent!55,arc=2mm,boxrule=0.5pt,left=2.5mm,right=2.5mm,top=1mm,bottom=1mm]\small{\color{accent}\faLightbulb[regular]}\ \textbf{Recorda:}\ #1\end{tcolorbox}}
\newcommand{\res}[1]{\boxed{#1}}\newcommand{\rest}[1]{\fbox{#1}}

\def\fmneed{11}% que el títol d'una secció no quede sol al final del full
% =====================================================================
% Capçalera FentMates de les fitxes: un color per a cada curs.
% Els colors són els de la pàgina inicial (LEVELS a assets/js/materials.js);
% el logo ve de _eines/logo/pdf/ (python3 _eines/logo/genera.py).
%
% Ús, al final del preàmbul de la unitat:
%   % =====================================================================
% Capçalera FentMates de les fitxes: un color per a cada curs.
% Els colors són els de la pàgina inicial (LEVELS a assets/js/materials.js);
% el logo ve de _eines/logo/pdf/ (python3 _eines/logo/genera.py).
%
% Ús, al final del preàmbul de la unitat:
%   % =====================================================================
% Capçalera FentMates de les fitxes: un color per a cada curs.
% Els colors són els de la pàgina inicial (LEVELS a assets/js/materials.js);
% el logo ve de _eines/logo/pdf/ (python3 _eines/logo/genera.py).
%
% Ús, al final del preàmbul de la unitat:
%   \input{../comu/fentmates.tex}
%   \fmcurs{1eso}
% i, al començament del document:
%   \capcalera{Matemàtiques}{1r ESO}{Fitxa d'exercicis · Unitat 2}{Divisibilitat}
% Opcional abans de \input: \def\fmneed{14} (espai mínim abans de cada secció).
% =====================================================================
\makeatletter
\def\fm@logodir{../../logo/pdf}
\providecommand{\fmneed}{10}
\def\fm@cursnom{}

% clau del curs, color i color del text damunt del color
\newcommand{\fm@curs}[3]{\@namedef{fm@color@#1}{#2}\@namedef{fm@text@#1}{#3}}
\fm@curs{1eso}{E6B84A}{navy}
\fm@curs{2eso}{3B82F6}{white}
\fm@curs{3eso}{FF7A50}{white}
\fm@curs{4eso-a}{A78BFA}{white}
\fm@curs{4eso-b}{6D28D9}{white}
\fm@curs{1batx-mi}{9B1C1C}{white}
\fm@curs{1batx-ccss}{EF6B6B}{white}
\fm@curs{2batx-mii}{1E6B47}{white}
\fm@curs{2batx-ccss}{2BB3A3}{white}

\newcommand{\fmcurs}[1]{%
  \@ifundefined{fm@color@#1}{\PackageError{fentmates}{Curs desconegut: #1}{}}{}%
  \edef\fm@tmp{\noexpand\definecolor{curs}{HTML}{\@nameuse{fm@color@#1}}%
               \noexpand\colorlet{curstext}{\@nameuse{fm@text@#1}}}\fm@tmp
  \colorlet{cursfosc}{curs!72!black}%
  \def\fm@logo{\fm@logodir/vertical-#1.pdf}%
  \def\fm@marca{\fm@logodir/marca-#1.pdf}}

% ---------- capçalera de la primera pàgina ----------
% #1 matèria (petit), #2 curs, #3 tipus de document i unitat, #4 títol
\newcommand{\capcalera}[4]{%
  \gdef\fm@cursnom{#2}%
  \par\noindent
  \begin{tcolorbox}[colback=curs,colframe=curs,boxrule=0pt,arc=2.2mm,width=3.5cm,height=2.15cm,
      valign=center,left=3mm,right=2mm,top=1mm,bottom=1mm,nobeforeafter,box align=center]
    \color{curstext}{\footnotesize #1}\par\vspace{1pt}{\fontsize{19}{21}\selectfont\bfseries #2\par}%
  \end{tcolorbox}\hspace{5mm}%
  \begin{minipage}[c]{\dimexpr\linewidth-3.5cm-5mm-2.9cm\relax}
    {\color{cursfosc}\small\bfseries\MakeUppercase{#3}}\par\vspace{2pt}
    {\color{navy}\fontsize{20}{23}\selectfont\bfseries #4\par}
  \end{minipage}\hfill
  \begin{minipage}[c]{2.6cm}\raggedleft\includegraphics[width=2.6cm]{\fm@logo}\end{minipage}\par
  \vspace{2pt}}

% ---------- barres de secció del color del curs ----------
\renewcommand{\seccio}[2]{\par\needspace{\fmneed\baselineskip}\bigskip
  \begin{tcolorbox}[colback=curs,colframe=curs,coltext=curstext,arc=2mm,boxrule=0pt,left=3mm,right=3mm,top=1.2mm,bottom=1.2mm]
  {\large\bfseries #1\quad #2}
  \end{tcolorbox}\vspace{-2pt}}

% ---------- estrelles de dificultat del color del curs (a Batxillerat no se'n posen) ----------
\@ifundefined{estrelles}{}{%
  \renewcommand{\estrelles}[1]{{\color{curs}\ifcase#1\or\faStar\faStar[regular]\faStar[regular]\or\faStar\faStar\faStar[regular]\or\faStar\faStar\faStar\fi}}}

% ---------- capçalera i peu de les pàgines ----------
% primera pàgina: logos institucionals al peu
\fancypagestyle{primera}{\fancyhf{}\renewcommand{\headrulewidth}{0pt}
  \fancyfoot[L]{\raisebox{-0.32\height}{\includegraphics[height=0.78cm]{logo-gva.png}}\hspace{5mm}%
    \raisebox{-0.32\height}{\includegraphics[height=0.78cm]{logo-vermellar.png}}}
  \fancyfoot[R]{\scriptsize\color{muted}\href{https://creativecommons.org/licenses/by-nc-sa/4.0/deed.ca}{\ccbyncsa}\ \ CC BY-NC-SA 4.0 · fentmates.org\qquad\small\thepage}}
% la resta: el símbol a l'esquerra i el curs a la dreta, amb el títol al centre com abans
\fancyhead[L]{\raisebox{-0.3\height}{\includegraphics[height=0.62cm]{\fm@marca}}}
\fancyhead[R]{\small\bfseries\color{cursfosc}\fm@cursnom}
\renewcommand{\headrule}{\color{curs!45}\hrule height \headrulewidth}
\makeatother

%   \fmcurs{1eso}
% i, al començament del document:
%   \capcalera{Matemàtiques}{1r ESO}{Fitxa d'exercicis · Unitat 2}{Divisibilitat}
% Opcional abans de \input: \def\fmneed{14} (espai mínim abans de cada secció).
% =====================================================================
\makeatletter
\def\fm@logodir{../../logo/pdf}
\providecommand{\fmneed}{10}
\def\fm@cursnom{}

% clau del curs, color i color del text damunt del color
\newcommand{\fm@curs}[3]{\@namedef{fm@color@#1}{#2}\@namedef{fm@text@#1}{#3}}
\fm@curs{1eso}{E6B84A}{navy}
\fm@curs{2eso}{3B82F6}{white}
\fm@curs{3eso}{FF7A50}{white}
\fm@curs{4eso-a}{A78BFA}{white}
\fm@curs{4eso-b}{6D28D9}{white}
\fm@curs{1batx-mi}{9B1C1C}{white}
\fm@curs{1batx-ccss}{EF6B6B}{white}
\fm@curs{2batx-mii}{1E6B47}{white}
\fm@curs{2batx-ccss}{2BB3A3}{white}

\newcommand{\fmcurs}[1]{%
  \@ifundefined{fm@color@#1}{\PackageError{fentmates}{Curs desconegut: #1}{}}{}%
  \edef\fm@tmp{\noexpand\definecolor{curs}{HTML}{\@nameuse{fm@color@#1}}%
               \noexpand\colorlet{curstext}{\@nameuse{fm@text@#1}}}\fm@tmp
  \colorlet{cursfosc}{curs!72!black}%
  \def\fm@logo{\fm@logodir/vertical-#1.pdf}%
  \def\fm@marca{\fm@logodir/marca-#1.pdf}}

% ---------- capçalera de la primera pàgina ----------
% #1 matèria (petit), #2 curs, #3 tipus de document i unitat, #4 títol
\newcommand{\capcalera}[4]{%
  \gdef\fm@cursnom{#2}%
  \par\noindent
  \begin{tcolorbox}[colback=curs,colframe=curs,boxrule=0pt,arc=2.2mm,width=3.5cm,height=2.15cm,
      valign=center,left=3mm,right=2mm,top=1mm,bottom=1mm,nobeforeafter,box align=center]
    \color{curstext}{\footnotesize #1}\par\vspace{1pt}{\fontsize{19}{21}\selectfont\bfseries #2\par}%
  \end{tcolorbox}\hspace{5mm}%
  \begin{minipage}[c]{\dimexpr\linewidth-3.5cm-5mm-2.9cm\relax}
    {\color{cursfosc}\small\bfseries\MakeUppercase{#3}}\par\vspace{2pt}
    {\color{navy}\fontsize{20}{23}\selectfont\bfseries #4\par}
  \end{minipage}\hfill
  \begin{minipage}[c]{2.6cm}\raggedleft\includegraphics[width=2.6cm]{\fm@logo}\end{minipage}\par
  \vspace{2pt}}

% ---------- barres de secció del color del curs ----------
\renewcommand{\seccio}[2]{\par\needspace{\fmneed\baselineskip}\bigskip
  \begin{tcolorbox}[colback=curs,colframe=curs,coltext=curstext,arc=2mm,boxrule=0pt,left=3mm,right=3mm,top=1.2mm,bottom=1.2mm]
  {\large\bfseries #1\quad #2}
  \end{tcolorbox}\vspace{-2pt}}

% ---------- estrelles de dificultat del color del curs (a Batxillerat no se'n posen) ----------
\@ifundefined{estrelles}{}{%
  \renewcommand{\estrelles}[1]{{\color{curs}\ifcase#1\or\faStar\faStar[regular]\faStar[regular]\or\faStar\faStar\faStar[regular]\or\faStar\faStar\faStar\fi}}}

% ---------- capçalera i peu de les pàgines ----------
% primera pàgina: logos institucionals al peu
\fancypagestyle{primera}{\fancyhf{}\renewcommand{\headrulewidth}{0pt}
  \fancyfoot[L]{\raisebox{-0.32\height}{\includegraphics[height=0.78cm]{logo-gva.png}}\hspace{5mm}%
    \raisebox{-0.32\height}{\includegraphics[height=0.78cm]{logo-vermellar.png}}}
  \fancyfoot[R]{\scriptsize\color{muted}\href{https://creativecommons.org/licenses/by-nc-sa/4.0/deed.ca}{\ccbyncsa}\ \ CC BY-NC-SA 4.0 · fentmates.org\qquad\small\thepage}}
% la resta: el símbol a l'esquerra i el curs a la dreta, amb el títol al centre com abans
\fancyhead[L]{\raisebox{-0.3\height}{\includegraphics[height=0.62cm]{\fm@marca}}}
\fancyhead[R]{\small\bfseries\color{cursfosc}\fm@cursnom}
\renewcommand{\headrule}{\color{curs!45}\hrule height \headrulewidth}
\makeatother

%   \fmcurs{1eso}
% i, al començament del document:
%   \capcalera{Matemàtiques}{1r ESO}{Fitxa d'exercicis · Unitat 2}{Divisibilitat}
% Opcional abans de \input: \def\fmneed{14} (espai mínim abans de cada secció).
% =====================================================================
\makeatletter
\def\fm@logodir{../../logo/pdf}
\providecommand{\fmneed}{10}
\def\fm@cursnom{}

% clau del curs, color i color del text damunt del color
\newcommand{\fm@curs}[3]{\@namedef{fm@color@#1}{#2}\@namedef{fm@text@#1}{#3}}
\fm@curs{1eso}{E6B84A}{navy}
\fm@curs{2eso}{3B82F6}{white}
\fm@curs{3eso}{FF7A50}{white}
\fm@curs{4eso-a}{A78BFA}{white}
\fm@curs{4eso-b}{6D28D9}{white}
\fm@curs{1batx-mi}{9B1C1C}{white}
\fm@curs{1batx-ccss}{EF6B6B}{white}
\fm@curs{2batx-mii}{1E6B47}{white}
\fm@curs{2batx-ccss}{2BB3A3}{white}

\newcommand{\fmcurs}[1]{%
  \@ifundefined{fm@color@#1}{\PackageError{fentmates}{Curs desconegut: #1}{}}{}%
  \edef\fm@tmp{\noexpand\definecolor{curs}{HTML}{\@nameuse{fm@color@#1}}%
               \noexpand\colorlet{curstext}{\@nameuse{fm@text@#1}}}\fm@tmp
  \colorlet{cursfosc}{curs!72!black}%
  \def\fm@logo{\fm@logodir/vertical-#1.pdf}%
  \def\fm@marca{\fm@logodir/marca-#1.pdf}}

% ---------- capçalera de la primera pàgina ----------
% #1 matèria (petit), #2 curs, #3 tipus de document i unitat, #4 títol
\newcommand{\capcalera}[4]{%
  \gdef\fm@cursnom{#2}%
  \par\noindent
  \begin{tcolorbox}[colback=curs,colframe=curs,boxrule=0pt,arc=2.2mm,width=3.5cm,height=2.15cm,
      valign=center,left=3mm,right=2mm,top=1mm,bottom=1mm,nobeforeafter,box align=center]
    \color{curstext}{\footnotesize #1}\par\vspace{1pt}{\fontsize{19}{21}\selectfont\bfseries #2\par}%
  \end{tcolorbox}\hspace{5mm}%
  \begin{minipage}[c]{\dimexpr\linewidth-3.5cm-5mm-2.9cm\relax}
    {\color{cursfosc}\small\bfseries\MakeUppercase{#3}}\par\vspace{2pt}
    {\color{navy}\fontsize{20}{23}\selectfont\bfseries #4\par}
  \end{minipage}\hfill
  \begin{minipage}[c]{2.6cm}\raggedleft\includegraphics[width=2.6cm]{\fm@logo}\end{minipage}\par
  \vspace{2pt}}

% ---------- barres de secció del color del curs ----------
\renewcommand{\seccio}[2]{\par\needspace{\fmneed\baselineskip}\bigskip
  \begin{tcolorbox}[colback=curs,colframe=curs,coltext=curstext,arc=2mm,boxrule=0pt,left=3mm,right=3mm,top=1.2mm,bottom=1.2mm]
  {\large\bfseries #1\quad #2}
  \end{tcolorbox}\vspace{-2pt}}

% ---------- estrelles de dificultat del color del curs (a Batxillerat no se'n posen) ----------
\@ifundefined{estrelles}{}{%
  \renewcommand{\estrelles}[1]{{\color{curs}\ifcase#1\or\faStar\faStar[regular]\faStar[regular]\or\faStar\faStar\faStar[regular]\or\faStar\faStar\faStar\fi}}}

% ---------- capçalera i peu de les pàgines ----------
% primera pàgina: logos institucionals al peu
\fancypagestyle{primera}{\fancyhf{}\renewcommand{\headrulewidth}{0pt}
  \fancyfoot[L]{\raisebox{-0.32\height}{\includegraphics[height=0.78cm]{logo-gva.png}}\hspace{5mm}%
    \raisebox{-0.32\height}{\includegraphics[height=0.78cm]{logo-vermellar.png}}}
  \fancyfoot[R]{\scriptsize\color{muted}\href{https://creativecommons.org/licenses/by-nc-sa/4.0/deed.ca}{\ccbyncsa}\ \ CC BY-NC-SA 4.0 · fentmates.org\qquad\small\thepage}}
% la resta: el símbol a l'esquerra i el curs a la dreta, amb el títol al centre com abans
\fancyhead[L]{\raisebox{-0.3\height}{\includegraphics[height=0.62cm]{\fm@marca}}}
\fancyhead[R]{\small\bfseries\color{cursfosc}\fm@cursnom}
\renewcommand{\headrule}{\color{curs!45}\hrule height \headrulewidth}
\makeatother

\fmcurs{1batx-mi}

\begin{document}
\thispagestyle{primera}

\capcalera{Matemàtiques I}{1r Batx}{Solucions de la fitxa · Unitat 8}{\raggedright Límits i continuïtat}
\vspace{8pt}

\begin{tcolorbox}[colback=soft,colframe=soft,arc=2mm,left=3mm,right=3mm,top=1.5mm,bottom=1.5mm]
\small Intenta resoldre cada exercici abans de mirar-ne la solució. Si no t'ix el mateix resultat, repassa els passos o comprova'l amb la calculadora de límits dels apunts interactius de la unitat.
\end{tcolorbox}

% =====================================================
\seccio{1}{Límits a l'infinit}

\exf{1} \begin{tasks}(2)
  \task $\displaystyle\lim_{x\to+\infty}=\res{+\infty}$, \ $\displaystyle\lim_{x\to-\infty}=\res{-\infty}$
  \task $\displaystyle\lim_{x\to+\infty}=\res{+\infty}$, \ $\displaystyle\lim_{x\to-\infty}=\res{-2}$
  \task $\displaystyle\lim_{x\to\pm\infty}=\res{2}$ \ (asímptota $y=2$)
  \task Oscil·la: \ $\res{\nexists}$ cap dels dos límits
\end{tasks}

\ex{2} \begin{tasks}(2)
  \task $-2x^3+7x^2$: \ $\res{+\infty}$ \ i \ $\res{-\infty}$
  \task $3x^4-7x$: \ $\res{+\infty}$ \ i \ $\res{+\infty}$
  \task $10^x$: \ $\res{0}$ \ i \ $\res{+\infty}$
  \task $\sqrt{5x-8}$: \ $\res{\nexists}$ \ i \ $\res{+\infty}$
  \task $\sqrt{-2x^2+1}$: \ $\res{\nexists}$ \ i \ $\res{\nexists}$
  \task $-5^x$: \ $\res{0}$ \ i \ $\res{-\infty}$
  \task $\sqrt{3-x}$: \ $\res{+\infty}$ \ i \ $\res{\nexists}$
  \task $\dfrac{x^2+3}{-x^3}$: \ $\res{0}$ \ i \ $\res{0}$
  \task $\dfrac{-x^3}{x^2+3}$: \ $\res{+\infty}$ \ i \ $\res{-\infty}$
  \task $\dfrac{5x^3-10}{3x^3+10x^2}$: \ $\res{\frac53}$ \ i \ $\res{\frac53}$
\end{tasks}\par\small (A la e), el domini és $\left[-\frac{1}{\sqrt2},\frac{1}{\sqrt2}\right]$; a la f), $-5^x=-(5^x)$.)

\ex{1} \begin{tasks}(3)
  \task $\lim\frac{x^2}{3x}=\lim\frac x3=\res{+\infty}$
  \task $\lim\frac{x^2}{x^3}=\res{0}$
  \task Mateix grau: $\res{\frac32}$
  \task $\lim\frac{4x^3}{-x^2}=\lim(-4x)=\res{+\infty}$
  \task $\frac{-6}{3}=\res{-2}$
  \task $\lim\frac{2x^2}{-x}=\lim(-2x)=\res{+\infty}$
\end{tasks}

\ex{2} \begin{tasks}(3)
  \task $\left(\tfrac13\right)^{+\infty}=\res{0}$
  \task $3^{-\infty}+1=\res{1}$
  \task $0{,}4^{-\infty}=\res{+\infty}$
  \task $\lim\frac{2x}{x}=\res{2}$
  \task $\sqrt{x^2}=|x|=-x$: \ $\lim\frac{-x}{2x}=\res{-\frac12}$
  \task $e^{-\infty}+2=\res{2}$
\end{tasks}


% =====================================================
\seccio{2}{Comparació d'infinits}

\exf{1} $\res{\ln x\ll\sqrt x\ll x^2\ll x^5\ll1{,}1^x\ll3^x}$

\ex{2} \begin{tasks}(3)
  \task $\res{0}$
  \task $\res{+\infty}$
  \task $\res{0}$
  \task $\res{-\infty}$
  \task $\lim\frac{2^x}{2\cdot2^x}=\res{\frac12}$
  \task $\res{+\infty}$
  \task $\res{-\infty}$
  \task $\lim\left(\frac54\right)^x=\res{+\infty}$
\end{tasks}

\ex{2} \begin{tasks}(2)
  \task $\sqrt{x^5}=x^{5/2}$ i $\frac52>2$: \ $\res{-\infty}$
  \task $3^x$ té major ordre: \ $\res{+\infty}$
\end{tasks}


% =====================================================
\seccio{3}{Indeterminacions a l'infinit}

\exf{2} \begin{tasks}(2)
  \task $\lim\frac{x+2x}{3x}=\res{1}$
  \task $\lim\frac{2x^2}{x^2}=\res{2}$
  \task $\lim\frac{3x}{x}=\res{3}$
  \task $\lim\frac{2x^{3/2}}{8x^2}=\lim\frac{1}{4\sqrt x}=\res{0}$
\end{tasks}

\ex{2} \begin{tasks}(1)
  \task $\dfrac{2x^2-5x-2x(x+3)}{x+3}=\dfrac{-11x}{x+3}\to\res{-11}$
  \task $\dfrac{x^2(x+2)-x^2(x-2)}{x^2-4}=\dfrac{4x^2}{x^2-4}\to\res{4}$
  \task $\dfrac{x^3-x(x^2-1)}{x^2-1}=\dfrac{x}{x^2-1}\to\res{0}$
  \task $\dfrac{(x^2+1)(3x-1)+x(5x-x^2)}{x(3x-1)}=\dfrac{2x^3+4x^2+3x-1}{3x^2-x}\to\lim\dfrac{2x}{3}=\res{+\infty}$
\end{tasks}

\ex{2} \begin{tasks}(1)
  \task $\dfrac{x^2-(x^2-1)}{x+\sqrt{x^2-1}}=\dfrac{1}{x+\sqrt{x^2-1}}\to\res{0}$
  \task $\dfrac{(x+2)-(x-2)}{\sqrt{x+2}+\sqrt{x-2}}=\dfrac{4}{\sqrt{x+2}+\sqrt{x-2}}\to\res{0}$
  \task $\dfrac{6x}{\sqrt{x^2+6x}+x}\to\dfrac{6x}{2x}=\res{3}$
  \task $\dfrac{-2-x}{\sqrt{x^2-2}+\sqrt{x^2+x}}\to\dfrac{-x}{2x}=\res{-\frac12}$
  \task $\dfrac{4x^2-(4x^2-3x)}{2x+\sqrt{4x^2-3x}}=\dfrac{3x}{2x+\sqrt{4x^2-3x}}\to\dfrac{3x}{4x}=\res{\frac34}$
\end{tasks}

\ex{2} \begin{tasks}(2)
  \task $\frac{x}{x+2}-1=\frac{-2}{x+2}$; \ $\lim\frac{-2(2x+1)}{x+2}=-4$: \ $\res{e^{-4}=\frac{1}{e^4}}$
  \task $\frac{-3}{2x+1}\cdot3x^2=\frac{-9x^2}{2x+1}\to-\infty$: \ $e^{-\infty}=\res{0}$
  \task $e^{\lim\frac3x\cdot2x}=\res{e^6}$
  \task $\frac{3}{x-2}\cdot x\to3$: \ $\res{e^3}$
  \task $\frac{4}{x^2-3}\cdot x^2\to4$: \ $\res{e^4}$
  \task La base tendeix a $3$: \ $3^{+\infty}=\res{+\infty}$ (no és $1^\infty$)
\end{tasks}

\ex{2} \begin{tasks}(1)
  \task $\displaystyle\lim_{x\to+\infty}\frac{12x+3\,000}{x}=12$: el cost per unitat s'acosta a \rest{12 €}.
  \task $\displaystyle\lim_{t\to+\infty}\frac{8t}{t^2+4}=0$: \rest{el medicament desapareix de la sang}.
  \task $P(0)=\res{300}$ conills. \ $\displaystyle\lim_{t\to+\infty}P(t)=500$: la població s'acosta a \rest{500 conills} sense superar-los ($P(t)=500-\frac{200}{t+1}$).
\end{tasks}


% =====================================================
\seccio{4}{Límits en un punt}

\exf{1} \label{ex:graf}\begin{tasks}(4)
  \task $\res{1}$
  \task $\res{+\infty}$
  \task $\res{-\infty}$
  \task $\res{-1}$
  \task $\res{2}$
  \task $\res{3}$
  \task $\res{\nexists}$
  \task $\res{1}$
\end{tasks}\par $\res{\nexists g(-2)}$, \ $g(1)=\res{2}$, \ $g(3)=\res{-1}$. \ També $\nexists\displaystyle\lim_{x\to-2}g(x)$ (laterals diferents).

\ex{1} $x=0{,}9;\ 0{,}99;\ 0{,}999$: \ $\frac{1}{x-1}=-10;\ -100;\ -1\,000\to\res{-\infty}$. \ \ $x=-1{,}1;\ -1{,}01;\ -1{,}001$: \ $x^2+5=6{,}21;\ 6{,}0201;\ 6{,}002001\to\res{6}$.

\ex{1} \begin{tasks}(3)
  \task $8-4+1=\res{5}$
  \task $\frac{4}{-2}=\res{-2}$
  \task $\frac60$: \ $-\infty$ per l’esquerra i $+\infty$ per la dreta: \ $\res{\nexists}$
  \task $\frac{-1}{0^+}$: \ $\res{-\infty}$ pels dos costats
  \task $\frac30$: \ $+\infty$ per l’esquerra i $-\infty$ per la dreta: \ $\res{\nexists}$
  \task $\frac{-1}{0^+}$: \ $\res{-\infty}$
  \task $\ln(0^+)=\res{-\infty}$
\end{tasks}

\ex{2} \begin{tasks}(1)
  \task $\lim_{x\to1}(2x+5)=\res{7}$ (el tros de $x<3$); \ $\lim_{x\to5}(-x+2)=\res{-3}$; \ en $x=3$: $\lim_{x\to3^-}(2x+5)=11$ i $\lim_{x\to3^+}(-x+2)=-1$: \ $\res{\nexists}$.
  \task $x=-1$: $\lim_{x\to-1^-}(x^2+1)=2$ i $\lim_{x\to-1^+}(1-x)=2$: \ $\res{2}$. \ \ $x=2$: $\lim_{x\to2^-}(1-x)=-1$ i $\lim_{x\to2^+}(x-3)=-1$: \ $\res{-1}$.
\end{tasks}


% =====================================================
\seccio{5}{Indeterminacions en un punt}

\exf{2} \begin{tasks}(2)
  \task $\frac{x-1}{(x-1)(x-2)}=\frac{1}{x-2}\to\res{-1}$
  \task $\frac{x(x-3)(x+3)}{x^2(x-3)}=\frac{x+3}{x}\to\res{2}$
  \task $\frac{(x-2)(x+2)}{(x-2)(x+3)}\to\res{\frac45}$
  \task $\frac{(x-1)(x^2+x+1)}{(x-1)(x+1)}\to\res{\frac32}$
  \task $\frac{x(x+2)}{(x+2)^2}=\frac{x}{x+2}\to\frac{-2}{0}$: \ $+\infty$ per l’esquerra, $-\infty$ per la dreta, $\res{\nexists}$
  \task $\frac{(x-1)^2(x+2)}{(x-1)^2(x+1)}=\frac{x+2}{x+1}\to\res{\frac32}$
\end{tasks}

\ex{2} \begin{tasks}(2)
  \task $\frac{x-9}{(x-9)(\sqrt x+3)}=\frac{1}{\sqrt x+3}\to\res{\frac16}$
  \task $\frac{x-4}{(x-4)(\sqrt{x+5}+3)}\to\res{\frac16}$
  \task $\frac{(x-1)(\sqrt{x+3}+2)}{x-1}=\sqrt{x+3}+2\to\res{4}$
  \task $\frac{4-(x+1)}{(x-3)(2+\sqrt{x+1})}=\frac{-1}{2+\sqrt{x+1}}\to\res{-\frac14}$
\end{tasks}

\ex{2} \begin{tasks}(2)
  \task $3x\cdot\frac2x=6$: \ $\res{e^6}$
  \task $(x-2)\cdot\frac{1}{x-2}=1$: \ $\res{e}$
  \task $\frac{x-1}{4}\cdot\frac{2}{x-1}=\frac12$: \ $\res{e^{1/2}=\sqrt e}$
  \task $(1-x)\cdot\frac{x}{x-1}=-x\to-1$: \ $\res{e^{-1}=\frac1e}$
\end{tasks}

\ex{2} En substituir: $\frac10-\frac20$, és a dir, $\infty-\infty$. \ Restem: $\frac{x+1-2}{x^2-1}=\frac{x-1}{(x-1)(x+1)}=\frac{1}{x+1}\to\res{\frac12}$.


% =====================================================
\seccio{6}{Continuïtat i tipus de discontinuïtat}

\exf{1} \begin{tasks}(4)
  \task \rest{Asimptòtica}
  \task \rest{Salt finit}
  \task \rest{Evitable} ($f(c)\neq\lim$)
  \task \rest{Evitable} ($\nexists f(c)$)
\end{tasks}

\ex{1} $x=-2$: un lateral és $-\infty$ i l'altre $-1$: \rest{salt infinit}. \ $x=1$: laterals $2$ i $3$: \rest{salt finit}. \ $x=3$: $\lim_{x\to3}g(x)=1\neq g(3)=-1$: \rest{evitable}. \ És contínua en la resta de punts.

\ex{1} \begin{tasks}(1)
  \task $\nexists f(0)$; \ $\lim_{x\to0}\frac{x(x-1)}{x}=-1$: \ \rest{evitable}
  \task $\nexists g(2)$; \ $\frac30$: $-\infty$ per l’esquerra i $+\infty$ per la dreta: \ \rest{asimptòtica}
  \task $h(0)=1$, \ $\lim_{x\to0^-}=1$, \ $\lim_{x\to0^+}=1$: \ \rest{contínua}
  \task $k(1)=3$, \ $\lim_{x\to1^-}(x+2)=3$, \ $\lim_{x\to1^+}\frac{1}{x-1}=+\infty$: \ \rest{salt infinit}
  \task $p(2)=3$, \ $\lim_{x\to2}(x+2)=4\neq3$: \ \rest{evitable}
\end{tasks}

\ex{2} \begin{tasks}(1)
  \task $x^2+x-2=0\Rightarrow x=1,\ x=-2$. \ $x=1$: $\frac{(x-1)(x+1)}{(x-1)(x+2)}\to\frac23$: \rest{evitable}; \ $x=-2$: $\frac{3}{0}$: \rest{asimptòtica}. Contínua en $\mathbb R-\{-2,1\}$.
  \task $x^2+4\neq0$: \ \rest{contínua en tot $\mathbb R$}.
  \task $x(x-4)=0\Rightarrow x=0,\ x=4$. \ $x=0$: $\frac{1}{x-4}\to-\frac14$: \rest{evitable}; \ $x=4$: $\frac{4}{0}$: \rest{asimptòtica}.
  \task Contínua on està definida: \ \rest{$[4,+\infty)$}.
\end{tasks}

\ex{2} \begin{tasks}(1)
  \task No està definida en $x=-5$, on $\frac{15}{0}$: \ \rest{asimptòtica en $x=-5$}; contínua en $\mathbb R-\{-5\}$.
  \task $g(x)=x+2$ si $x\neq0$ i no està definida en $0$: \ \rest{evitable en $x=0$}.
  \task $\lim_{x\to3^-}h=1\neq\lim_{x\to3^+}h=4$: \ \rest{salt finit en $x=3$}; contínua en $\mathbb R-\{3\}$.
  \task $x=0$: $1=1=k(0)$; \ $x=2$: $\lim_{x\to2^-}(1-x^2)=-3=k(2)$: \ \rest{contínua en tot $\mathbb R$}.
\end{tasks}


% =====================================================
\seccio{7}{Funcions a trossos. Paràmetres}

\exf{2} \begin{tasks}(1)
  \task Trossos polinòmics: contínues. \ $x=2$: $f(2)=-1$, $\lim_{x\to2^-}(x-3)=-1$, $\lim_{x\to2^+}(x^2-5)=-1$: \ \rest{contínua en $\mathbb R$}.
  \task El denominador $x-1$ s’anul·la en $1$, fora de $x\geq2$. \ $x=0$: $1=1=g(0)$, contínua. \ $x=2$: $\lim_{x\to2^-}=5$, $\lim_{x\to2^+}=3=g(2)$: \ \rest{salt finit en $x=2$}.
  \task $\frac1x$ no està definida en $x=0$, que és del seu interval: $-\infty$ i $+\infty$, \rest{asimptòtica en $x=0$}. \ $x=-1$: $h(-1)=1$, laterals $1$ i $-1$: \rest{salt finit}. \ $x=1$: $h(1)=1$, laterals $1$ i $1$: contínua.
\end{tasks}

\ex{2} $x^2-4<0$ si $-2<x<2$. \ Per a $x>2$: $f(x)=\frac{(x-2)(x+2)}{x-2}=x+2\to4$; \ per a $-2<x<2$: $f(x)=-(x+2)\to-4$. \ $\nexists f(2)$ i els laterals són $-4$ i $4$: \ \rest{salt finit en $x=2$}.

\ex{2} $\lim_{x\to2^-}(x^2-3)=1=f(2)$; \ $\lim_{x\to2^+}(-x^2+m)=-4+m$. \ $-4+m=1\Rightarrow\res{m=5}$.

\ex{2} \begin{tasks}(1)
  \task $a+3=1-a\Rightarrow\res{a=-1}$
  \task $\lim_{x\to3}(x+3)=6\Rightarrow\res{k=6}$
  \task $1+a=\frac42=2\Rightarrow\res{a=1}$ \ (el denominador $x+2$ s’anul·la en $-2$, fora de $x>0$)
\end{tasks}

\ex{3} $x=0$: $a=1$. \ $x=2$: $2b+1=4-a=3\Rightarrow b=1$. \ $\res{a=1,\ b=1}$

\ex{2} $\lim_{x\to10^-}P(x)=3+15=18$ i $\lim_{x\to10^+}P(x)=10a+8$: \ $10a+8=18\Rightarrow\res{a=1}$ €/km. \ $P(25)=25+8=\res{33\text{ €}}$.


% =====================================================
\seccio{8}{Branques infinites i asímptotes}

\exf{2} \begin{tasks}(1)
  \task Cap AV ($x^2+4\neq0$). \ AH $\res{y=1}$: \ $f(x)-1=\frac{-5}{x^2+4}<0$, la corba queda \rest{per davall}. No té AO.
  \task AV $\res{x=1}$: $\frac{1}{0^+}$, $+\infty$ pels dos costats. \ AH $\res{y=0}$: per damunt si $x\to+\infty$ ($g>0$) i per davall si $x\to-\infty$. No té AO.
  \task AV $\res{x=3}$: $-\infty$ per l’esquerra i $+\infty$ per la dreta. \ No té AH. \ $2x^2=(x-3)(2x+6)+18$: \ AO $\res{y=2x+6}$; \ $\frac{18}{x-3}$: per damunt si $x\to+\infty$ i per davall si $x\to-\infty$.
\end{tasks}

\ex{2} \begin{tasks}(1)
  \task AV $\res{x=2}$ ($-\infty$ per l’esquerra, $+\infty$ per la dreta); AH $\res{y=3}$ ($f(x)-3=\frac{7}{x-2}$: per damunt a la dreta i per davall a l’esquerra).
  \task AV $\res{x=-2}$ ($+\infty$ per l’esquerra, $-\infty$ per la dreta) i $\res{x=2}$ ($-\infty$ per l’esquerra, $+\infty$ per la dreta); AH $\res{y=1}$.
  \task En $x=-1$: $\frac00$, $\frac{1}{x-1}\to-\frac12$: no hi ha AV (forat). \ AV $\res{x=1}$ ($-\infty$ per l’esquerra, $+\infty$ per la dreta); AH $\res{y=0}$.
  \task AV $\res{x=1}$ ($-\infty$ per l’esquerra, $+\infty$ per la dreta). \ $x^2+x+1=(x-1)(x+2)+3$: AO $\res{y=x+2}$.
  \task AV $\res{x=0}$ ($\frac{1}{0^+}$: $+\infty$ pels dos costats). \ $f(x)=x+\frac{1}{x^2}$: AO $\res{y=x}$, i la corba queda sempre per damunt.
  \task No és racional: \ $\lim_{x\to+\infty}\frac{3\cdot2^x}{2^x}=3$ i $\lim_{x\to-\infty}\frac{3\cdot0}{0+1}=0$: \ AH $\res{y=3}$ per la dreta i $\res{y=0}$ per l’esquerra. No té AV.
\end{tasks}

\ex{3} $x\to+\infty$: $m=\lim\frac{\sqrt{x^2+4x}}{x}=1$, \ $n=\lim(\sqrt{x^2+4x}-x)=\lim\frac{4x}{\sqrt{x^2+4x}+x}=2$: \ $\res{y=x+2}$. \ \ $x\to-\infty$: $\sqrt{x^2}=-x$, $m=-1$, \ $n=\lim(\sqrt{x^2+4x}+x)=\lim\frac{4x}{\sqrt{x^2+4x}-x}=\frac{4x}{-2x}=-2$: \ $\res{y=-x-2}$.

\ex{2} \begin{tasks}(1)
  \task $\lim f=+\infty$, $m=\lim(x^2-3)=+\infty$: \ \rest{branca parabòlica}
  \task $\lim g=+\infty$, $m=\lim\frac{\ln x}{x}=0$: \ \rest{branca parabòlica}
  \task $\lim h=+\infty$, $m=\lim\frac{x^2}{x-1}=+\infty$: \ \rest{branca parabòlica}
  \task $k(x)=x+\frac1x$: \ \rest{asímptota obliqua $y=x$}, no és branca parabòlica
\end{tasks}

\ex{3} AV $x=-2$: el denominador s'anul·la en $-2$: $c=2$. \ AH: $\displaystyle\lim_{x\to\pm\infty}\frac{ax+b}{x+2}=a=3$. \ Punt: $f(0)=\frac b2=1\Rightarrow b=2$. \ $\res{a=3,\ b=2,\ c=2}$


\end{document}
